\section{Implementation Notes and Analysis Protocols}
\label{app:hyper}

This appendix records the temperature ladder and the graph, walk, and analysis details deferred from Section~\ref{sec:method}. The implementation has no separate teacher and student temperature knobs for transition-row or ambient matching: each such term reuses the temperature with which its target was constructed.

\subsection{The rest of the temperature ladder}
\label{app:hyper-ladder}

The $r=1$ readout uses the graph bandwidth, $\tau_1(i)=\tau_i$. Broader scales are set by $\tau_r(i)=\sqrt{r}\,\tau_i$, following the $\sqrt{t}$ spread of a diffusion at nominal time $r$. This is a rule rather than a derivation: the lazy walk takes $r/2$ real steps in expectation, so a strict derivation would carry a different constant. The ladder therefore has no free entries and automatically inherits entropic per-row bandwidths. Scale weights follow the complementary rule $\omega_r=1/r$ with $\omega_0=\omega_1$, so influence falls with diffusion time.

\subsection{Truncation tolerance}
\label{app:hyper-truncation}

Sparse diffusion, the per-anchor candidate pool, and the walk support all perform the same operation: keep a subset of a probability row and renormalize. RIPPLE fixes one tolerance $\delta$ and keeps, per row, the smallest column set carrying $1-\delta$ of the mass. Such a truncation perturbs the row by exactly $\delta$ in total variation and by $-\log(1-\delta)$ nats in Kullback--Leibler divergence---a bound in the units of the loss itself---and by Lemma~\ref{lem:truncation} the error is additive across steps, so an $r$-step target is within $r\delta$ of the exact one. Every stored array width is therefore a consequence of $\delta$. Memory guards cap the widest rows; the build reports if either bound binds before the tolerance is met, in which case the guarantee does not hold.

\subsection{Non-backtracking walks}
\label{app:hyper-nonbacktracking}

The walk of Eq.~\eqref{eq:walk} is non-backtracking: the step at $t+1$ is drawn from $P^T_{j_t}$ conditioned on $j_{t+1}\neq j_{t-1}$, except at a node of degree one, where the edge just traversed is the only one available. Since rows are supervised with weight equal to their visit count, an excursion that returns immediately spends the walk budget re-supervising a row it has already seen; forbidding it raises the number of distinct rows a walk of given length reaches, and non-backtracking walks also have faster mixing and cover time on graphs of bounded degree~\citep{alon2007nonbacktracking}. The constraint changes only the occupancy measure $\nu_B$---which rows are chosen---and never a target: a visited node $j$ is still matched against its own transition row $P^T_j$ at $\tau_j$. Proposition~\ref{prop:walkgauge} is likewise unaffected, since it counts \emph{traversed edges} and a repeated edge contributes no new constraint. 

\subsection{Dense rather than sampled walk targets}
\label{app:hyper-dense}

A natural alternative scores the sampled transitions themselves, $-\log p^{S}(j_{t+1}\mid j_{t})$ summed along each path. The walk is first-order and each step is scored on the teacher's own $j_t$, so that path likelihood factorizes into one-step terms whose expectation over the walk is exactly Eq.~\eqref{eq:lrow}: the sampled form is an unbiased but higher-variance estimator of the dense one. Since $P^{T}_{j}$ is already materialized during graph construction, we use the dense form and pay no variance. Nothing in the term couples two consecutive steps; it matches one-step transition rows at walk-visited nodes.

\subsection{Scale-separation diagnostic}
\label{app:scale-separation}

On a common candidate pool for every scale, the diagnostic in Table~\ref{tab:multiscale_ablation} is
\begin{equation}
E_{\mathrm{sep}}
=\frac{1}{|\mathcal A|\binom{|\mathcal R|}{2}}
\sum_{i\in\mathcal A}\sum_{r<s}
\left|
\mathrm{JS}(p^T_{i,r},p^T_{i,s})
-\mathrm{JS}(q^S_{i,r},q^S_{i,s})
\right|,
\qquad \mathcal R=\{1,2,4\}.
\end{equation}
Using a common pool is essential: otherwise changing support could create apparent separation even when two student readouts assign identical scores to shared columns.

\subsection{Per-setting row-loss grid}
\label{app:lambda-grid}

Table~\ref{tab:lambda_sweep_all} expands the aggregate in Table~\ref{tab:lambda_sweep}. Only $\lambda$ changes and every cell remains the downstream nine-task average; no calibration diagnostic is used for model selection.

\begin{table}[H]
\centering
\caption{Per-setting row-loss sweep. Entries are downstream Avg. (mean $\pm$ standard deviation). Results are pending.}
\label{tab:lambda_sweep_all}
\setlength{\tabcolsep}{7pt}
\begin{tabular}{lccccc}
\toprule
Setting $\backslash\ \lambda$ & $0$ & $0.1$ & $0.3$ & $0.5$ & $1.0$ \\
\midrule
Qwen-4B $\rightarrow$ BERT          & \tbd & \tbd & \tbd & \tbd & \tbd \\
BGE-M3 $\rightarrow$ MiniLM-768     & \tbd & \tbd & \tbd & \tbd & \tbd \\
Qwen-0.6B $\rightarrow$ MiniLM-384 & \tbd & \tbd & \tbd & \tbd & \tbd \\
\bottomrule
\end{tabular}
\end{table}

\subsection{Sampling-efficiency measurement protocol}
\label{app:sampling-efficiency-protocol}

For a fixed probe batch, define the gradient fidelity of an $m$-candidate objective as
\begin{equation}
\label{eq:gradient-fidelity}
G_m=\cos\!\left(\nabla_\theta\mathcal L_{\mathrm{diff}}^{(m)},
\nabla_\theta\mathcal L_{\mathrm{diff}}^{(\mathrm{full})}\right).
\end{equation}
Only $m$ changes across Table~\ref{tab:m_sweep}; the sampler and all other training choices are fixed. Target mass is measured before renormalizing the sampled row. $G_m$ is computed at the same checkpoint and on the same probe batch for every $m$; the full-support gradient is used only for evaluation. Step time excludes graph preprocessing and is averaged after warm-up, while peak accelerator memory is reset immediately before the measured steps. All timing and memory results use the same hardware and software stack.

\begin{table}[H]
\centering
\caption{Candidate-budget mechanism test. Only $m$ changes. Results are pending.}
\label{tab:m_sweep}
\setlength{\tabcolsep}{4pt}
\begin{tabular}{rccccc}
\toprule
$m$ & Target mass$\uparrow$ & $G_m\uparrow$ & Avg.$\uparrow$ & Time/step$\downarrow$ & Memory$\downarrow$ \\
\midrule
16  & \tbd & \tbd & \tbd & \tbd & \tbd \\
32  & \tbd & \tbd & \tbd & \tbd & \tbd \\
64  & \tbd & \tbd & \tbd & \tbd & \tbd \\
128 & \tbd & \tbd & \tbd & \tbd & \tbd \\
256 & \tbd & \tbd & \tbd & \tbd & \tbd \\
\bottomrule
\end{tabular}
\end{table}
